% Compilar con pdfLaTeX dos veces para resolver el índice y las referencias.
% Documento autocontenido: no requiere imágenes externas ni archivo .bib.
\documentclass[11pt,a4paper]{article}

\usepackage[utf8]{inputenc}
\usepackage[T1]{fontenc}
\usepackage[spanish,es-tabla,es-nodecimaldot,es-lcroman]{babel}
\usepackage{lmodern}
\usepackage[margin=2.3cm,headheight=15pt]{geometry}
\usepackage{amsmath,amssymb}
\usepackage{array,tabularx}
\usepackage{color}
\usepackage{fancyhdr}
\usepackage{hyperref}

\definecolor{azul}{rgb}{0.08,0.20,0.34}
\definecolor{grisclaro}{rgb}{0.94,0.96,0.98}
\definecolor{gristexto}{rgb}{0.35,0.39,0.43}
\hypersetup{
  colorlinks=true,
  linkcolor=azul,
  citecolor=azul,
  urlcolor=azul,
  pdftitle={Metodología de visión por computadora y navegación con TurtleBot 4 Lite},
  pdfsubject={Informe metodológico del proyecto de Visión por Computadora}
}
\urlstyle{same}
\linespread{1.06}
\setlength{\parindent}{0pt}
\setlength{\parskip}{5pt}
\setlength{\emergencystretch}{2em}
\setcounter{tocdepth}{2}
\setcounter{totalnumber}{4}
\renewcommand{\topfraction}{0.95}
\renewcommand{\bottomfraction}{0.85}
\renewcommand{\textfraction}{0.05}
\renewcommand{\floatpagefraction}{0.85}
\setlength{\textfloatsep}{12pt}
\setlength{\intextsep}{10pt}
\clubpenalty=10000
\widowpenalty=10000
\raggedbottom
\renewcommand{\arraystretch}{1.22}
\newcolumntype{Y}{>{\raggedright\arraybackslash}X}
\newcolumntype{P}[1]{>{\raggedright\arraybackslash}p{#1}}
\renewcommand{\refname}{Referencias}

\pagestyle{fancy}
\fancyhf{}
\fancyhead[L]{\small Visión por Computadora}
\fancyhead[R]{\small TurtleBot 4 Lite}
\fancyfoot[R]{\footnotesize\thepage}
\renewcommand{\headrulewidth}{0.4pt}
\renewcommand{\footrulewidth}{0pt}

\newcommand{\pendiente}[1]{%
  \par\smallskip
  \begingroup
  \setlength{\fboxsep}{8pt}%
  \noindent\colorbox{grisclaro}{%
    \parbox{\dimexpr\linewidth-2\fboxsep\relax}{%
      \textbf{Por completar con resultados de competición.} #1}}%
  \par\endgroup\smallskip
}
\newcommand{\modulo}[2]{%
  \fbox{\parbox[c][2.15cm][c]{\dimexpr0.27\textwidth-2\fboxsep-2\fboxrule\relax}{%
    \centering\small\textbf{#1}\par\smallskip #2}}%
}

\begin{document}

\hypersetup{pageanchor=false}
\begin{titlepage}
  \thispagestyle{empty}
  \centering
  \vspace*{1.0cm}
  {\large Universidad de Ingeniería y Tecnología (UTEC)\par}
  \vspace{0.3cm}
  {\normalsize [Carrera / Escuela profesional]\par}
  \vspace{1.5cm}
  {\color{azul}\large VISIÓN POR COMPUTADORA\par}
  \vspace{0.5cm}
  {\color{azul}\rule{\textwidth}{1pt}\par}
  \vspace{0.65cm}
  {\fontsize{22}{28}\selectfont\bfseries
    Navegación autónoma, control gestual\\
    y reconocimiento facial con\\
    TurtleBot 4 Lite\par}
  \vspace{0.65cm}
  {\color{azul}\rule{\textwidth}{1pt}\par}
  \vspace{1.1cm}
  {\large Informe metodológico del proyecto final\par}
  \vspace{0.35cm}
  {\normalsize Actualizado al 17 de noviembre de 2026\par}
  \vspace{1.0cm}
  \begin{tabular}{ll}
    \textbf{Integrantes:} & [Nombre y código del integrante 1]\\
                         & [Nombre y código del integrante 2]\\
                         & [Nombre y código del integrante 3]\\
                         & [Nombre y código del integrante 4]\\[0.6em]
    \textbf{Docente:} & Rensso Victor Hugo Mora Colque\\[0.6em]
    \textbf{Entrega:} & [Fecha anunciada en clase]
  \end{tabular}
  \vfill
  \begin{minipage}{0.90\textwidth}
    \small\color{gristexto}
    Este informe presenta la metodología, arquitectura, algoritmos de visión
    por computadora, estrategias de navegación y decisión, justificación de
    diseño y análisis de limitaciones del sistema desarrollado sobre TurtleBot~4
    Lite. Los detalles de implementación y trazabilidad del código fuente se
    ubican en el apéndice.
  \end{minipage}
  \vspace*{0.4cm}
\end{titlepage}

\pagenumbering{roman}
\hypersetup{pageanchor=true}
\section*{Resumen}
\addcontentsline{toc}{section}{Resumen}

Este documento describe la metodología de visión por computadora y control
robótico desarrollada sobre un TurtleBot~4 Lite con ROS~2 Jazzy para resolver
dos modalidades en circuitos con obstáculos revelados el día de la
competición~\cite{enunciado}: navegación autónoma (\textbf{Task~A}) y
teleoperación visual mediante gestos humanos (\textbf{Task~B}).

La solución se organiza en una arquitectura híbrida desacoplada que separa el
procesamiento visual en CPU (laptop externa) de la capa de actuación, fusión
sensorial y seguridad en tiempo real sobre ROS~2. Para el reconocimiento facial
se emplean \textbf{YuNet} y \textbf{SFace} con normalización $\ell_2$,
similitud coseno y criterio calibrado de umbral y margen con rechazo explícito.
Para el control por gestos se aplica \textbf{MediaPipe Hand Landmarker} con
clasificación geométrica en coordenadas de píxeles, histéresis de
$0.4\ \mathrm{s}$ y caducidad monotónica de $0.5\ \mathrm{s}$. Para la
navegación autónoma se integra evasión reactiva conservadora por sectores
LiDAR (exigiendo cobertura mínima del $30\%$), detección visual de señales
direccionales (bordes Canny, segmentación HSV por tramo y análisis geométrico
de flechas) y maniobras de giro de $90^\circ$ controladas por odometría. Se
justifican las decisiones de ingeniería frente a alternativas evaluadas, se
resumen los resultados experimentales obtenidos hasta la fecha y se discuten
las limitaciones actuales junto con las líneas de mejora futura.

\textbf{Palabras clave:} visión por computadora, metodología, TurtleBot 4 Lite,
ROS~2, reconocimiento facial, reconocimiento de gestos, navegación autónoma.

\tableofcontents
\clearpage
\pagenumbering{arabic}

\section{Introducción y planteamiento del problema}

La operación de un robot móvil en circuitos desconocidos mediante visión por
computadora requiere transformar flujo visual y lecturas métricas en decisiones
consistentes de percepción y movimiento bajo restricciones de latencia,
iluminación variable y seguridad física. De acuerdo con los lineamientos del
proyecto~\cite{enunciado}, la configuración del circuito y las fotografías de
evaluación se dan a conocer recién en la competición, por lo que la metodología
no puede basarse en trayectorias memorizadas ni en coordenadas fijas.

El sistema debe resolver dos modalidades operativas que comparten el subsistema
de identificación facial:
\begin{itemize}
  \item \textbf{Task~A (navegación autónoma):} recibir el nombre de una persona
  objetivo perteneciente a la galería antes de iniciar; registrar el punto de
  partida~A; recorrer de forma autónoma un circuito con obstáculos utilizando
  flechas azules como guía direccional hasta confirmar la meta~B; localizar y
  confirmar la fotografía de la persona solicitada entre varias candidatas;
  leer únicamente el código QR asociado a dicha fotografía; y retornar guiado
  por flechas rojas hasta confirmar~A para mostrar la palabra obtenida.
  \item \textbf{Task~B (teleoperación por gestos):} completar el recorrido
  $A \rightarrow B \rightarrow \text{identificación facial en } B \rightarrow A
  \rightarrow \text{indicación del nombre}$ conduciendo al robot únicamente con
  señales visuales humanas capturadas en video, sin controles manuales ni
  maniobras autónomas durante el intento.
\end{itemize}

\section{Arquitectura general de la solución}
\label{sec:arquitectura}

El sistema sigue una \textbf{arquitectura híbrida modular} que separa la
percepción y toma de decisiones de alto nivel (ejecutadas en CPU en la laptop
del operador) de la ejecución de bajo nivel sobre ROS~2 Jazzy (ejecutada en un
contenedor aislado conectado por red al TurtleBot~4 Lite). Esta separación
responde a dos necesidades metodológicas: poder validar los algoritmos de
visión de manera reproducible sin acceso continuo al robot físico y aislar la
seguridad de actuación mediante una única autoridad de velocidad.

\begin{figure}[htbp]
  \centering
  \setlength{\tabcolsep}{4pt}
  \begin{tabular}{@{}ccccc@{}}
    \multicolumn{5}{c}{\small\textbf{Modalidad Task A: Navegación Autónoma y Reconocimiento}}\\
    \noalign{\vskip 4pt}
    \modulo{Percepción visual}{Cámara del robot\\YuNet + SFace\\Detector de flechas}
    & $\longrightarrow$ &
    \modulo{Supervisor de misión}{Objetivo nominal\\Estados e histéresis\\Validación de trama}
    & $\longrightarrow$ &
    \modulo{Control y sensores}{Puerta única de salida\\Evasión LiDAR + TF\\Giros por odometría}\\[10pt]
    \noalign{\vskip 4pt}
    \multicolumn{5}{c}{\small\textbf{Modalidad Task B: Teleoperación Gestual y Actuación Física}}\\
    \noalign{\vskip 4pt}
    \modulo{Entrada gestual}{Webcam del operador\\MediaPipe (21 puntos)\\Reglas y filtro 0.4 s}
    & $\longrightarrow$ &
    \modulo{Puente de seguridad}{Vigilante 20 Hz\\Caducidad 0.5 s\\Reloj monotónico}
    & $\longrightarrow$ &
    \modulo{TurtleBot 4 Lite}{Comandos de velocidad\\Base Create~3 + LiDAR\\Cámara OAK-D}
  \end{tabular}
  \caption{Arquitectura general del sistema desacoplada por planos de percepción,
  decisión y actuación. Elaboración propia.}
  \label{fig:arquitectura}
\end{figure}

La comunicación entre el host y el entorno ROS~2 se realiza mediante flujos
estructurados con marcas de tiempo monotónicas compartidas, manteniendo la
recepción de video de la cámara del robot en un canal independiente limitado a
$5\ \mathrm{FPS}$ para visualización y análisis sin interferir con el envío de
comandos de parada. La Tabla~\ref{tab:componentes} resume la responsabilidad y
el estado de verificación de cada subsistema.

\begin{table}[htbp]
  \centering
  \small
  \begin{tabularx}{\textwidth}{@{}P{4.2cm}P{5.4cm}Y@{}}
    \hline
    \textbf{Subsistema} & \textbf{Función metodológica} & \textbf{Estado de validación}\\
    \hline
    Reconocimiento facial &
    Detección, alineación, descripción y decisión con rechazo explícito &
    Validado localmente con galería ampliada y pruebas unitarias\\
    Entrada nominal y panel &
    Selección previa del objetivo y confirmación en vivo en 3 fotogramas &
    Probado en ventana visual y ensayos combinados en robot\\
    Percepción gestual (Task~B) &
    Estimación de mano, reglas geométricas y confirmación temporal &
    Confirmado en simulación Gazebo y en hardware real\\
    Puente ROS y caducidad &
    Traducción a comandos con estampa de tiempo y parada por expiración &
    Confirmado en simulación y pruebas físicas preliminares\\
    Evasión LiDAR y odometría &
    Análisis de cobertura por sectores, TF y control angular de giros &
    Probado en pruebas unitarias y movimiento físico preliminar\\
    Detección de señales de giro &
    Identificación de flechas azules y rojas con prioridad ante muro &
    Validado en 17 pruebas unitarias y láminas de referencia\\
    Supervisor integral Task~A &
    Detección de meta A/B por etapa, barrido en B, lector QR y vuelta &
    En integración final sobre la base modular actual\\
    \hline
  \end{tabularx}
  \caption{Subsistemas de la arquitectura propuesta y estado de madurez.}
  \label{tab:componentes}
\end{table}

\section{Algoritmos de visión por computadora}

\subsection{Reconocimiento facial: YuNet, SFace y criterio dual con rechazo}
\label{sec:algoritmo_facial}

El subsistema facial sigue un flujo de cuatro etapas diseñado para operar en
CPU con colecciones de tamaño reducido sin reentrenar redes profundas~\cite{opencv}:

\begin{enumerate}
  \item \textbf{Preprocesamiento y detección:} cada imagen se reescala en
  memoria manteniendo su relación de aspecto hasta un lado máximo de
  $960\ \mathrm{px}$. El detector \textbf{YuNet} (arquitectura convolucional
  ligera para detección de rostros y cinco puntos clave: ojos, nariz y
  comisuras) localiza rostros candidatos exigiendo una confianza mínima de
  detección y un tamaño facial mínimo en píxeles. Para evitar ambigüedades, una
  consulta solo es válida si contiene \emph{exactamente un rostro útil}; cero o
  múltiples rostros producen rechazo inmediato.

  \item \textbf{Alineación afín y extracción de características:} usando los
  cinco puntos faciales de YuNet se aplica alineación geométrica y recorte
  normalizado, sobre el cual \textbf{SFace} extrae un vector descriptor de
  $128$ dimensiones $f(x)\in\mathbb{R}^{128}$. Cada descriptor se proyecta a la
  hiperesfera unitaria mediante normalización $\ell_2$:
  \begin{equation}
    z(x)=\frac{f(x)}{\lVert f(x)\rVert_2},
    \qquad c(x,r)=z(x)^\top z(r),
    \label{eq:coseno}
  \end{equation}
  donde $c(x,r)\in[-1,1]$ mide la similitud coseno entre la consulta $x$ y una
  referencia $r$ de la galería.

  \item \textbf{Agregación por identidad:} para aumentar la robustez frente a
  variaciones de pose o expresión sin diluir la puntuación por muestras menos
  favorables, cada identidad $i$ se puntúa con la media de sus $k_i=\min(2,n_i)$
  referencias más similares en el conjunto $T_i(x)$:
  \begin{equation}
    S_i(x)=\frac{1}{k_i}\sum_{r\in T_i(x)}c(x,r).
  \end{equation}

  \item \textbf{Regla de decisión con margen y rechazo:} sea $i^\star=\arg\max_i
  S_i(x)$ el mejor candidato y $\Delta=S_{i^\star}(x)-\max_{j\ne i^\star}S_j(x)$
  su margen frente al segundo mejor. Se acepta $i^\star$ únicamente si supera
  simultáneamente un umbral absoluto de similitud $\tau$ y un margen mínimo de
  separación $\mu$:
  \begin{equation}
    \widehat y =
    \begin{cases}
      i^\star, & S_{i^\star}(x)\geq\tau\ \text{y}\ \Delta\geq\mu,\\
      \text{NO\_CONCLUYENTE}, & \text{en otro caso}.
    \end{cases}
    \label{eq:decision}
  \end{equation}
  Los parámetros $(\tau,\mu)$ se calibran exclusivamente sobre un subconjunto de
  calibración independiente de la galería y de la evaluación reservada,
  priorizando cero falsas aceptaciones (incluidos negativos simulados al
  retirar temporalmente la identidad correcta). En operación en vivo sobre el
  robot, la confirmación de una identidad exige coincidencia sostenida en tres
  fotogramas nuevos consecutivos.
\end{enumerate}

\subsection{Reconocimiento de gestos: MediaPipe y clasificación geométrica}
\label{sec:algoritmo_gestos}

Para Task~B se emplea \textbf{MediaPipe Hand Landmarker}~\cite{mediapipe}, que
estima en dos etapas la región de la palma y las coordenadas tridimensionales
de $21$ puntos clave de la mano. Sobre dichos puntos se diseñó un clasificador
geométrico interpretable en lugar de entrenar una red clasificadora adicional:

\begin{itemize}
  \item \textbf{Normalización geométrica a píxeles:} las coordenadas
  normalizadas $[0,1]^2$ entregadas por MediaPipe se escalan al ancho y alto
  reales de la captura antes de evaluar distancias y ángulos, evitando sesgos
  por relación de aspecto rectangular.
  \item \textbf{Reglas de extensión y dirección:} se evalúa el estado de
  extensión de los cuatro dedos largos comparando las articulaciones
  metacarpofalángicas y las puntas. El índice extendido hacia arriba con los
  otros tres dedos flexionados ordena \textbf{Avanzar}; el índice extendido
  lateralmente ordena \textbf{Girar} en el sentido correspondiente al operador
  frente a la cámara (invirtiendo una única vez la lectura horizontal en la
  interpretación sin aplicar espejo visual a la imagen); la palma abierta, la
  ausencia de mano, más de una mano o configuraciones en banda angular ambigua
  producen \textbf{Detener}.
  \item \textbf{Estabilidad temporal:} toda orden de movimiento requiere
  mantenerse estable durante $0.4\ \mathrm{s}$ para activarse; cualquier cambio
  de sentido obliga a pasar por parada y reiniciar la confirmación, mientras que
  la palma abierta detiene de forma instantánea.
\end{itemize}

\subsection{Percepción visual de señales direccionales para Task A}

Para interpretar los carteles circulares de giro del circuito (flechas azules en
el tramo de ida y flechas rojas en el de vuelta), se implementó una tubería
clásica de visión geométrica libre de entrenamiento supervisado:

\begin{enumerate}
  \item \textbf{Detección de candidatos circulares:} la imagen BGR de la cámara
  OAK-D se reescala a un ancho máximo de $640\ \mathrm{px}$ y se extraen bordes
  Canny sobre una combinación máxima entre luminancia y canales azul/rojo para
  resaltar ambos colores de señal frente al fondo. Se ajustan elipses
  candidatas con diámetro mínimo de $55\ \mathrm{px}$ y se descartan carteles
  recortados contra los bordes del encuadre.
  \item \textbf{Verificación cromática por tramo:} dentro de cada elipse se
  muestrea el espacio HSV (separando azul en $H\in[90,135]$ y rojo en las dos
  bandas extremas de matiz $H\le 12$ o $H\ge 168$) y se descartan discos que no
  correspondan al color activo del tramo actual (\texttt{azul} en ida,
  \texttt{rojo} en vuelta).
  \item \textbf{Análisis geométrico de flecha:} dentro de la región interior se
  evalúa la simetría y orientación de segmentos y vértices de punta para
  clasificar \texttt{IZQUIERDA} o \texttt{DERECHA}, rechazando discos vacíos,
  carteles de prohibición/STOP y texturas de fondo.
  \item \textbf{Confirmación temporal y rearme:} una maniobra solo se dispara si
  un único cartel del color correcto se detecta de forma consistente en tres
  fotogramas nuevos consecutivos, y no se rearma hasta acumular al menos
  $1.5\ \mathrm{s}$ de ausencia visual y $0.4\ \mathrm{m}$ de desplazamiento.
\end{enumerate}

\section{Estrategias de navegación autónoma, sensores y toma de decisiones}

\subsection{Integración sensorial y autoridad única}

La seguridad del robot se apoya en cuatro flujos sensoriales integrados bajo el
principio de \textbf{autoridad única de velocidad}: un único nodo supervisor en
ROS~2 publica mensajes \texttt{TwistStamped} en \texttt{/cmd\_vel}, anulando el
movimiento a $(0,0)$ ante la expiración o invalidez de cualquier fuente:

\begin{itemize}
  \item \textbf{RPLIDAR (\texttt{/scan}):} transformado al marco del chasis
  (\texttt{base\_link} $\rightarrow$ \texttt{rplidar\_link}) mediante el árbol
  estático de transformaciones (TF) para corregir la rotación de montaje del
  sensor. Toda lectura con antigüedad $>0.6\ \mathrm{s}$ bloquea el avance.
  \item \textbf{Odometría (\texttt{/odom}):} validada en finitud numérica,
  norma unitaria de cuaternión y estabilidad de marcos de referencia (un cambio
  inesperado de marco revoca el permiso de movimiento). Registra la posición
  inicial~A tras confirmar sensores válidos y mide la variación angular en giros.
  \item \textbf{Cámara OAK-D:} entregada con perfil QoS de sensores en canal
  separado, validando dimensiones de trama y caducidad de recepción de
  $1.0\ \mathrm{s}$.
  \item \textbf{Latido de control del host:} el proceso de percepción envía
  órdenes de estado con caducidad de $0.5\ \mathrm{s}$ medidas sobre reloj
  monotónico compartido; la pérdida de latido o el fin de flujo detiene el
  robot.
\end{itemize}

\subsection{Evasión reactiva conservadora y control de giros (Task A)}

La estrategia de navegación autónoma combina evasión local basada en LiDAR con
maniobras discretas guiadas por las señales visuales y cerradas por odometría:

\begin{enumerate}
  \item \textbf{Evasión LiDAR por cobertura de sectores:} el escaneo se divide
  en un cono frontal ($\pm 35^\circ$) y dos sectores laterales. A diferencia de
  enfoques ingenuos que consideran cualquier rayo sin retorno ($+\infty$) como
  espacio libre, cada sector exige una \textbf{cobertura mínima del $30\%$} de
  lecturas finitas (al menos $3$ rayos útiles) y $\le 10\%$ de errores
  (\texttt{NaN} o fuera de rango). Si un sector carece de cobertura suficiente,
  se declara no observable y fuerza parada inmediata. Con cobertura válida, el
  robot avanza a $0.10\ \mathrm{m/s}$, reduce o corrige ante precaución a
  $0.45\ \mathrm{m}$ y detiene el avance si un obstáculo frontal entra a
  $0.32\ \mathrm{m}$ o un lateral a $0.28\ \mathrm{m}$, girando a
  $0.40\ \mathrm{rad/s}$ hacia el costado más despejado.

  \item \textbf{Arbitraje con señales direccionales:} cuando aparece una flecha
  confirmada correspondiente al color del tramo, su dirección tiene prioridad
  sobre la preferencia lateral de la evasión ante un muro. El robot ejecuta un
  giro controlado de $90^\circ \pm 5^\circ$ integrando la diferencia angular de
  odometría (con manejo continuo de cruces en $\pm\pi$). Si durante el giro la
  holgura lateral o frontal cae por debajo de $0.32\ \mathrm{m}$ o falla el
  LiDAR/TF, se detiene sin invertir jamás el sentido ordenado por la señal; si
  excede la tolerancia angular o agota $12\ \mathrm{s}$, enclava parada de
  seguridad.

  \item \textbf{Transición hacia reconocimiento y retorno:} durante el prototipo
  actual, al detectar el rostro objetivo en tres fotogramas nuevos se pausa de
  inmediato el desplazamiento. Para completar el contrato integral de
  competición, esta lógica se extiende con la detección de señal de meta para
  separar el recorrido en etapas (evitando confundir A y B), un barrido
  rotacional en B para ubicar el rostro solicitado entre varias fotos, la
  lectura exclusiva del QR vinculado a esa fotografía y la conmutación al tramo
  de señales rojas para regresar a~A.
\end{enumerate}

\section{Justificación de modelos, algoritmos y decisiones de diseño}

La Tabla~\ref{tab:justificacion} resume los principales criterios técnicos que
guiaron la selección de modelos y los compromisos de ingeniería asumidos frente
a otras alternativas.

\begin{table}[htbp]
  \centering
  \small
  \begin{tabularx}{\textwidth}{@{}P{3.4cm}P{4.2cm}Y@{}}
    \hline
    \textbf{Elección adoptada} & \textbf{Alternativa considerada} & \textbf{Justificación técnica y compromiso}\\
    \hline
    YuNet + SFace preentrenados en CPU &
    Entrenamiento propio o inferencia en nube/GPU &
    Opera sin conexión a internet (requisito del laboratorio) ni GPU, con
    alineación facial integrada; exige cuidar iluminación y calidad de galería.\\
    Decisión dual $(\tau,\mu)$ con rechazo &
    Vecino más cercano que siempre etiqueta &
    Minimiza penalizaciones por falsa identificación en competición, aceptando
    rechazar fotogramas borrosos o ambiguos.\\
    MediaPipe + reglas geométricas (Task~B) &
    Clasificador neuronal de gestos propio &
    Evita recolectar y etiquetar miles de manos, ofrece interpretabilidad directa
    y permite imponer bandas muertas de seguridad y confirmación de $0.4\ \mathrm{s}$.\\
    Visión geométrica Canny/HSV para flechas &
    Detector de objetos YOLO reentrenado &
    No requiere dataset de entrenamiento de carteles y permite ajustar
    tolerancias de color por tramo; resulta sensible a variaciones extremas de
    iluminación o oclusión parcial.\\
    Navegación reactiva + señales + odometría frente a solo SLAM/Nav2 &
    Dependencia única de \texttt{slam\_toolbox} + Nav2 &
    Aunque se evaluó una base SLAM/Nav2 con exploración por fronteras, su
    ajuste en pasillos estrechos con desfases de reloj Create~3/Pi resultó
    frágil y no incorporaba flechas ni QR; la estrategia reactiva da control
    directo, aunque exige integrar explícitamente la meta y el retorno.\\
    Cobertura LiDAR mínima del $30\%$ como bloqueo &
    Tratar rayos $+\infty$ como camino despejado &
    Evita colisiones ante zonas ciegas o pérdida de retornos en hardware real;
    redujo falsos paros al calibrarse de $50\%$ a $30\%$ tras medir $37$--$47\%$
    de retornos finitos.\\
    Procesamiento en host + canal de video a $5\ \mathrm{FPS}$ &
    Ejecución completa a bordo o video crudo sin límite &
    Facilita depuración y protege \texttt{/cmd\_vel} de congestión Wi-Fi;
    depende de mantener reloj monotónico compartido y conexión inalámbrica.\\
    \hline
  \end{tabularx}
  \caption{Justificación de selección de modelos, algoritmos y decisiones de diseño.}
  \label{tab:justificacion}
\end{table}

\section{Resultados experimentales preliminares}
\label{sec:resultados}

\subsection{Evaluación del reconocimiento facial}

La metodología facial se evaluó en cortes progresivos manteniendo siempre
separadas las referencias de galería, el conjunto de calibración y las
consultas reservadas (Tabla~\ref{tab:resultados_facial}).

\begin{table}[htbp]
  \centering
  \small
  \begin{tabularx}{\textwidth}{@{}P{4.5cm}cccY@{}}
    \hline
    \textbf{Corte experimental} & \textbf{Consultas} & \textbf{Aciertos} & \textbf{Errores} & \textbf{Rechazos y negativos simulados}\\
    \hline
    Colección inicial (26 fotos, 5 identidades) &
    5 & 4 & 0 &
    1 rechazo por similitud baja; 5/5 negativos simulados rechazados
    ($\tau=0.5726$, $\mu=0.2691$).\\
    Corte intermedio (6 identidades) &
    6 & 4 & 0 &
    2 rechazos (ambigüedad / similitud insuficiente); 6/6 negativos simulados
    rechazados ($\tau=0.5735$, $\mu=0.2498$).\\
    Galería ampliada (74 fotos, 13 identidades) &
    13 & 13 & 0 &
    0 rechazos; 13/13 negativos simulados rechazados (47 referencias válidas,
    13 calibración, 13 evaluación).\\
    \hline
  \end{tabularx}
  \caption{Resultados exploratorios de reconocimiento facial en particiones
  locales reservadas.}
  \label{tab:resultados_facial}
\end{table}

Los cortes muestran cero identificaciones incorrectas en todas las evaluaciones
y rechazo completo de los negativos simulados; sin embargo, al provenir de
colecciones locales pequeñas, aún deben validarse sistemáticamente con
fotografías impresas vistas desde la cámara del robot en pista y con personas
desconocidas independientes.

\subsection{Pruebas de integración, simulación y robot físico}

De manera paralela a la evaluación facial, el sistema cuenta con pruebas
automatizadas y ensayos físicos parciales:
\begin{itemize}
  \item \textbf{Verificación de software:} se registraron \textbf{164 pruebas
  unitarias} superadas en desarrollo (113 correspondientes a Task~A y 51 a
  Task~B), incluidas 17 pruebas específicas de detección y control de señales
  direccionales sobre las 10 láminas del archivo de referencia y sus reflejos
  horizontales.
  \item \textbf{Task~B (gestos):} se confirmó el control completo en simulación
  Gazebo y en movimiento físico sobre el robot real con envío de comandos
  \texttt{TwistStamped}, además de recepción de video comprimido en panel
  lateral a $5\ \mathrm{FPS}$.
  \item \textbf{Task~A (autónoma):} se confirmó movimiento físico en el robot
  real combinando evasión LiDAR y reconocimiento facial en vivo, ensayo en el
  cual se observaron coberturas finitas de $37$--$47\%$ y pausas puntuales de
  recepción de hasta $1.66\ \mathrm{s}$, lo que permitió ajustar los umbrales de
  seguridad y cobertura. Las pruebas físicas de giros con carteles impresos,
  lectura QR y retorno completo permanecen en ejecución de laboratorio.
\end{itemize}

\pendiente{Incorporar la tabla oficial de resultados en competición para
Task~A y Task~B: llegada a B, persona solicitada, palabra QR (Task~A), retorno a
A, tiempo total, colisiones, intervenciones y puntuación final.}

\section{Limitaciones y mejoras futuras}

A partir de las pruebas realizadas se identificaron las principales limitaciones
del sistema actual y sus líneas concretas de mejora:

\begin{enumerate}
  \item \textbf{Dependencia de iluminación y plano único del LiDAR:} la
  detección de flechas en HSV puede degradarse bajo sombras fuertes o reflejos, y
  el LiDAR bidimensional no percibe obstáculos muy bajos por debajo de su altura
  de montaje. \emph{Mejora futura:} incorporar calibración adaptativa de balance
  de blancos/luminancia y fusionar la nube de puntos de profundidad de la cámara
  OAK-D para obstáculos bajos.
  \item \textbf{Choques diagonales y huella del robot:} la evaluación por
  sectores discretos (frontal e izquierdo/derecho) no certifica toda la
  envolvente circular del chasis en giros cercanos a esquinas. \emph{Mejora
  futura:} proyectar la huella geométrica completa del robot sobre los puntos
  del escáner o reintegrar un mapa de costos local con control de colisiones.
  \item \textbf{Completitud de la secuencia final de Task~A:} el prototipo
  reconoce rostros durante el avance y todavía no cierra el ciclo completo de
  meta~B, barrido facial multi-foto, lectura QR asociada y vuelta autónoma a~A.
  \emph{Mejora futura:} integrar la máquina de estados por etapas con detector de
  meta, decodificador QR espacialmente acotado a la caja facial confirmada y
  navegación de retorno por señales rojas más verificación odométrica/visual
  de~A.
  \item \textbf{Latencia inalámbrica y desacoplamiento de red:} el transporte de
  video y órdenes entre laptop y robot depende de la estabilidad del router de
  laboratorio, con pausas observadas de hasta $1.66\ \mathrm{s}$ en recepción.
  \emph{Mejora futura:} medir cuantitativamente los tiempos de frenado y latencia
  de extremo a extremo, y migrar la inferencia liviana (detección de señales y
  parada de emergencia) directamente a bordo de la Raspberry~Pi del robot.
  \item \textbf{Generalización facial en impresiones:} las pruebas actuales se
  apoyan en fotografías digitales de colección pequeña y negativos simulados.
  \emph{Mejora futura:} ampliar la evaluación con carteles impresos a distintas
  distancias y ángulos capturados directamente por la cámara OAK-D del robot.
\end{enumerate}

\clearpage
\phantomsection
\addcontentsline{toc}{section}{Referencias}
\begin{thebibliography}{9}

\bibitem{enunciado}
R.~V.~H. Mora Colque.
\emph{Attention is not all you need for robots: Final project}.
Enunciado del curso, 13 de octubre de 2026.
Apartados \emph{Tasks Description}, \emph{Submission} y \emph{Rubric}.

\bibitem{opencv}
OpenCV.
\emph{DNN-based Face Detection and Recognition}.
\href{https://docs.opencv.org/4.x/d0/dd4/tutorial_dnn_face.html}
{Documentación oficial de detección, alineación y reconocimiento facial}.
Consulta: noviembre de 2026.

\bibitem{mediapipe}
Google AI Edge.
\emph{Hand Landmarker}.
\href{https://developers.google.com/edge/mediapipe/solutions/vision/hand_landmarker}
{Documentación oficial de estimación de puntos de la mano}.
Consulta: noviembre de 2026.

\bibitem{ros2}
ROS 2.
\emph{Topics y arquitectura DDS}, documentación Jazzy.
\href{https://github.com/ros2/ros2_documentation/blob/jazzy/source/Concepts/Basic/About-Topics.rst}
{Referencia oficial de comunicación mediante tópicos}.
Consulta: noviembre de 2026.

\end{thebibliography}

\clearpage
\appendix
\section{Detalles de implementación y trazabilidad de módulos}

En concordancia con el enfoque metodológico del cuerpo principal, este apéndice
resume la correspondencia entre los subsistemas descritos y los archivos de
implementación del repositorio. Los datos personales, pesos preentrenados
(\texttt{.onnx}, \texttt{.task}) y registros locales de sesión permanecen fuera
del control de versiones conforme a \texttt{.gitignore}.

\begin{table}[htbp]
  \centering
  \footnotesize
  \begin{tabularx}{\textwidth}{@{}P{6.2cm}Y@{}}
    \hline
    \textbf{Módulos de implementación} & \textbf{Detalle técnico y protocolos}\\
    \hline
    \texttt{task 1/reconocimiento\_\allowbreak facial.py}\\
    \texttt{task 1/modelos\_\allowbreak rostros.py}\\
    \texttt{task 1/01\_\allowbreak reconocimiento\_\allowbreak rostros.ipynb} &
    Carga verificada por SHA-256 de YuNet y SFace; partición auditable en
    \texttt{particion\_\allowbreak rostros.json}; calibración y evaluación
    separadas.\\
    \texttt{task 1/interfaz\_\allowbreak facial.py}\\
    \texttt{task 1/prueba\_\allowbreak facial\_\allowbreak robot.py} &
    Panel OpenCV en vivo con entrada de nombre objetivo sin tildes, sugerencias
    de galería y confirmación en tres fotogramas nuevos.\\
    \texttt{task 1/evasion\_\allowbreak lidar.py}\\
    \texttt{task 1/busqueda\_\allowbreak lidar\_\allowbreak ros.py}\\
    \texttt{task 1/control\_\allowbreak busqueda.py} &
    Lógica de cobertura LiDAR ($\ge 30\%$, máx. $10\%$ errores), verificación de
    TF (\texttt{base\_\allowbreak link} $\rightarrow$ \texttt{rplidar\_\allowbreak link}),
    validación de \texttt{/odom} y puerta única sobre \texttt{/cmd\_\allowbreak vel}
    (\texttt{TwistStamped}).\\
    \texttt{task 1/geometria\_\allowbreak senales.py}\\
    \texttt{task 1/senales\_\allowbreak giro.py}\\
    \texttt{task 1/control\_\allowbreak giros.py} &
    Detección geométrica de discos/flechas (HSV azul/rojo, Canny, elipses) y
    control odométrico de giros de $90^\circ \pm 5^\circ$ con enclavamiento de
    seguridad y rearme temporal/espacial.\\
    \texttt{task 1/busqueda\_\allowbreak robot.py}\\
    \texttt{task 1/ejecutar\_\allowbreak autonomo\_\allowbreak robot.sh} &
    Orquestador host--Docker con modos aislados (\texttt{--solo-evasion},
    \texttt{--solo-reconocimiento}, \texttt{--solo-senales},
    \texttt{--evasion-reconocimiento}) sin copiar datos faciales al contenedor.\\
    \texttt{task 2/gestos.py}\\
    \texttt{task 2/control\_\allowbreak gestos.py}\\
    \texttt{task 2/puente\_\allowbreak gestos\_\allowbreak ros.py}\\
    \texttt{task 2/webcam\_\allowbreak gestos.py} &
    Clasificador MediaPipe (21 puntos), filtro de $0.4\ \mathrm{s}$, puente
    JSON por \texttt{stdin} con vigilante monotónico a $20\ \mathrm{Hz}$ y
    caducidad de $0.5\ \mathrm{s}$.\\
    \texttt{task 2/camara\_\allowbreak robot.py}\\
    \texttt{task 2/camara\_\allowbreak robot\_\allowbreak ros.py} &
    Transporte separado de fotogramas OAK-D comprimidos hacia el host a un
    máximo de $5\ \mathrm{FPS}$ y caducidad de $1.0\ \mathrm{s}$.\\
    \texttt{task 1/tests/} y \texttt{task 2/tests/} &
    Conjunto de 164 pruebas automatizadas (lógica facial, señales, cobertura
    LiDAR, protocolos de tiempo real, gestos y transporte de cámara).\\
    \hline
  \end{tabularx}
  \caption{Correspondencia entre subsistemas metodológicos y artefactos de
  implementación.}
  \label{tab:apendice_archivos}
\end{table}

\section{Protocolo de comunicación y parámetros operativos}

Para garantizar la reproducibilidad en entorno de laboratorio sin conexión a
internet, la capa de red ROS~2 utiliza descubrimiento DDS en subred con par
estático explícito (\texttt{rmw\_\allowbreak fastrtps\_\allowbreak cpp},
\texttt{ROS\_\allowbreak AUTOMATIC\_\allowbreak DISCOVERY\_\allowbreak RANGE=SUBNET}
y \texttt{ROS\_\allowbreak STATIC\_\allowbreak PEERS}), sincronización horaria de
la Raspberry~Pi al inicio de cada sesión y contenedores en modo de red
\texttt{host}. Los parámetros activos de navegación autónoma y señales se
centralizan en archivos de configuración JSON (\texttt{evasion\_\allowbreak lidar.json}
y \texttt{senales\_\allowbreak giro.json}) para ajustar umbrales sin alterar el
código de control.

\end{document}
